\documentclass[11pt,a4paper]{article}
\usepackage[margin=25mm]{geometry}
\usepackage{amsmath,amssymb,amsthm,hyperref}
\setlength{\emergencystretch}{2em}
\newtheorem{proposition}{Proposition}
\newcommand{\gen}[1]{\langle#1\rangle}
\title{A singleton omission orbit at an interior crossing}
\author{Human Compression project: bounded L044 proof component}
\date{}
\begin{document}\maketitle
Let $G=\gen{\rho_0,\ldots,\rho_{r-1}}\leq A_n$ be a faithful
transitive string group generated by involutions, $n\geq8$,
with every maximal omission intransitive. Suppose the unique
eligible $i$-edge $\{a,b\}$ has $0<i<r-1$, and the two
$G_i$-orbits are $A$ and the singleton $\{b\}$.
Write $L=G_{<i}$, $R=G_{>i}$ and $d=r-1-i$.
All generators except $\rho_i$ fix $b$; those at distance
more than one from $i$ also fix $a$, by the unique-crossing
lemma. No IP for a quotient action will be used.

\begin{proposition}
Under these hypotheses, $r\leq\lfloor(n-1)/2\rfloor$.
\end{proposition}
\begin{proof}
Neither $L$ nor $R$ is transitive on $A$. Indeed, if $L$ were,
then $\gen{L,\rho_i}$ would be transitive on $A\cup\{b\}$,
contradicting intransitivity of any omission with index greater
than $i$. The argument for $R$ is the dual one. Both factors
are nontrivial, because the interior index leaves generators
on both sides and a redundant generator would have a transitive
omission.

Since $L$ and $R$ commute and $G_i=LR$ is transitive on $A$,
the $L$-orbits form $m\geq2$ blocks of common size $k\geq2$.
The action of $L$ on any one of these blocks is faithful: its
kernel fixes all $R$-translates of that block, hence all of $A$
and also $b$. Thus its $i$ distinguished generators retain
independence in degree $k$. Whiston gives $k\geq i+1$.

Every distinguished $R$ block image is essential. Otherwise
the remaining $R$ images, together with $L$, would still be
transitive on $A$. The omitted generator is supported on $A$
and so cannot join any orbits of that global omission. Adjoining
it produces transitive $G$, so the omission would already be
transitive, a contradiction.

In particular the last block image is nontrivial. The distance
to a moved-block argument gives $m\geq d+1$: start from
$\beta=a^L$, and use a shortest path of earlier $R$ block
images to a block moved by the last generator. Every intervening
index must occur, by commuting the factors on opposite sides
of a missing index; the upper factors fix $\beta$ and the
lower factors commute with the last generator. The path has
at least $d-1$ edges, and its last edge under $\rho_{r-1}$
gives a further block. For $d=1$ the same count is just
$m\geq2$. Consequently
\[
 n=km+1\geq(i+1)(d+1)+1
          =2r+(i-1)(d-1)-1.                             \tag{1}
\]

If $i,d\geq2$ and $(i-1)(d-1)\geq2$, (1) gives $n\geq2r+1$.
The remaining case with $i,d\geq2$ is $i=d=2$, so $r=5$.
Rank failure would force $n\leq10$, while (1) gives $n\geq10$;
hence $k=m=3$ and $n=10$. But each of the two distinguished
involutions of $L$ is a transposition on each faithful
three-point orbit. It is then a product of three transpositions
and fixes $b$, contrary to evenness.

We may therefore assume one of $i,d$ is one. Reverse the string
if necessary, so $i=1$. Then $L=\gen{\rho_0}$ has two-point
orbits throughout $A$: $k=2$, $n=2m+1$, and $\rho_0$ exchanges
all $m$ pairs. Its evenness forces $m$ even. In particular
$n\geq8$ gives $m\geq4$.
If the rank bound fails, $r=d+2>m$, while $m\geq d+1$.
Thus $d=m-1$ and $r=m+1$.

The essential $R$ block images form an independent family of
$m-1$ involutions in $S_m$, so their image is $S_m$, by
Whiston's equality clause. As already proved, $R$ is not
transitive on $A$. Its transitive action on the $m$ pairs then
forces exactly two point orbits $U,V$, each meeting each pair
once. The block map on $R$ is faithful, since a kernel element
fixes the one point in each pair belonging to each of $U,V$.
Thus $R$ is the diagonal natural $S_m$ on these two rows.

Let $a'=a^{\rho_0}$ and put $H=\gen{\rho_3,\ldots,\rho_{r-1}}$.
It fixes $a$ by unique crossing, and $a'$ by commutation with
$\rho_0$. Its $m-2$ block images remain independent, fix the
pair $\{a,a'\}$, and hence generate the full $S_{m-1}$ on the
remaining pairs, again by Whiston. In particular $H$ has only
the two fixed points $a,a'$ in $A$, and is diagonal natural
$S_{m-1}$ on the remaining two rows.

The crossing involution is $s=\rho_1=(a\ b)u$, where $u$ is
supported on $A\setminus\{a\}$. Since $H$ fixes $a,b$ and
commutes with $s$, the permutation $u$ centralizes $H$.
It fixes $a$, so it fixes $a'$ as well. On the remaining two
rows, each of size $m-1\geq3$, the centralizer of the diagonal
natural $S_{m-1}$ consists of the identity and the unique
equivariant row exchange. Indeed a commuting permutation
permutes its two orbits, and the natural symmetric action
has trivial centralizer. The exchange here is the restriction
of $\rho_0$, so is the product $v$ of the other $m-1$ pair
transpositions. Thus $u\in\{1,v\}$. Since $s$ is even,
$u=v$, which is odd because $m$ is even.

But $s$ now joins the two $R$-orbits $U,V$ through any one of
these other pairs, and joins $b$ to $a$. Consequently
$G_0=\gen{R,s}$ is transitive, the final contradiction.
This proves the bound in every case.
\end{proof}

\paragraph{Scope and attribution.}
This proves the interior singleton-orbit branch used from
Cameron--Fernandes--Leemans--Mixer, \emph{Highest rank of a
polytope for $A_n$}, arXiv:1605.09173v1, Proposition~5.25,
physical page~33. The source's commuting-block reduction and
the retained Whiston bound are reused; the actual-orbit count
and diagonal-row omission argument replace its path-disjointness
and equality-graph steps. The unique-crossing input is paid
in full. The moved-block argument is reproduced here so no
nontrivial-second-orbit hypothesis is imported from the other
support note. No maximal-string classification, finite census,
or projected IP is needed. The endpoint indices $i=0,r-1$
are explicitly outside this statement. There is no full L044,
new question, formalization or priority claim.
\end{document}
